% Integrals and their differentials: a differential is a finished factor, not a big operator, and an
% integral's operand spans the whole product after it.
\documentclass{article}
\usepackage{amsmath,amssymb}
\begin{document}

% formerly repro math-parse/differential_is_no_big_operator
% witness: 57al A/B Perl sample (root-causer): 1,596 formulas in 188 papers lost a differential — 2605.12296 S3.Ex68,
%          2605.21644 S4.Ex10, 2605.24070 S5.Ex112, 2605.26800 S3.E42Xb; `diffop_apply` gives the `d` of `d x` the DIFFOP
%          role, which `is_bigop_operator` (pragmatics.rs) counts as a big operator, so 57al's wider absorption pragma
%          pruned every product with a `d x` before another factor and the survivor read the `d` as a variable.
%          Perl's `IntFactor` `diffd` form (MathGrammar:643-646) is a finished factor; only DIFFOP lexemes (`\partial`,
%          iopart's `\rmd`) are big operators. The same role check made `\int_0^1 dx\,f` read `d * x * f`
%          (`is_bigop_or_scripted_bigop`, semantics.rs; older)
% oracle:  Perl 0 errors: `…@(g * differential-d@(u) * differential-d@(v))`, `integral@(P * x * differential-d@(mu) * x) = K`,
%          `…@(f * x * differential-d@(x) * g * y)`, `((integral _ 0) ^ 1)@(differential-d@(x) * f)`
% engines: rust 57al: `g * d * u * differential-d@(v)`, `P@(x) * d * mu * x`, `f@(x) * d * x * g * y`, `d * x * f`
%          (RUST-ONLY; `P@(x)` for Perl's `P * x` is divergence #18)
% expect:  0 errors; every differential kept, as Perl
% Perl `IntFactor` `diffd` form (MathGrammar:643-646) is a finished factor; only DIFFOP lexemes (`\partial`)
% are big operators. Every differential is kept; `P@(x)` for Perl's `P * x` is divergence #18.
\section{A differential is no big operator}
\[ \int_{0}^{1}\int_{0}^{1}g\,du\,dv \]
\[ \int P(x)\,d\mu(x)=K \]
\[ \int_{0}^{1}f(x)\,dx\,g(y) \]
\[ \int_0^1 dx\,f \]

% formerly repro math-parse/bigop_operand_spans_an_application_after_a_mulop (57ak probes, the PDE weak form)
% witness: 57ak probes (`\int_\Omega\nabla u\cdot\nabla v\,dx`, the weak form of PDE papers): a big operator's
%          operand stops before an application that follows a MulOp — `\int u\cdot\sin v` is
%          integral@(u) cdot sine@(v) — where Perl's `addOpArgs` / `addIntOpArgs` (MathGrammar:605-630) take
%          the whole term, a MulOp chain of Factors, applications included. Older than 57ak (57aj the same)
% oracle:  Perl 0 errors: `integral@(u cdot sine@(v))`, `integral@(nabla@(u) cdot nabla@(v))`,
%          `(sum _ i)@(a _ i cdot logarithm@(b _ i))`, `integral@(u cdot v)`; 57al review, a longer product:
%          `integral@(f cdot g cdot h)`, `a cdot (sum _ i)@(b _ i cdot c)`, `integral@(u cdot sine@(v cdot w))`
% expect:  0 errors; Perl's text for each
% Perl `addOpArgs` / `addIntOpArgs` (MathGrammar:605-630; `moreOpArgFactors` :612-617, `moreIntOpArgFactors` :633-638)
% take the whole term, applications included: `integral@(u cdot sine@(v))`, `a cdot (sum _ i)@(b _ i cdot c)`.
% Perl: `moreOpArgFactors` (MathGrammar:612-617; `moreIntOpArgFactors` :633-638) takes every factor after a MulOp
\section{A big operator's operand spans an application after a multiplication}
$\int u\cdot\sin v$ $\int\nabla u\cdot\nabla v$ $\sum_i a_i\cdot\log b_i$ $\int u\cdot v$
$\int f\cdot g\cdot h$ $a\cdot\sum_i b_i\cdot c$ $\int u\cdot\sin v\cdot w$

% formerly repro math-parse/gathered_row_keeps_its_differential (57x review; 47 of the 3,003 A/B papers)
% witness: 57x review; the content branch of gathered/split/multline (an XMWrap of XMRefs to the rows' tokens)
%          keeps an XMRef to `\,dx`'s `d` token after the row's parse replaced that token by
%          differential-d@(x); Perl's per-formula LOSTNODES remap (MathParser.pm:287-297) rewrites such refs, and
%          it is still a comment in latexml_math_parser/src/parser.rs (parse_formula); the final XML drops the `d`
%          from the content (`times(f, x)`) — every such row with a differential, since before 57x
% witnesses: 2605.01547, 2605.23309
% oracle:  pdflatex 0 errors
% engines: rust 57w3 and 57x `integral@(f * [] * x) = d`; perl `integral@(f * differential-d@(x)) = d` (RUST-ONLY)
% expect:  0 errors, 0 warnings; text="integral@(f * differential-d@(x)) = d"
% Perl parses the content branch too (MathParser.pm:378-392); its LOSTNODES remap (:287-297) and `Annotate`
% (:1206-1235) keep the XMRef to the `d`: "integral@(f * differential-d@(x)) = d".
\section{A gathered row keeps its differential}
\[\begin{gathered}\int f\,dx=d\end{gathered}\]

% witness: 2605.03741, 2605.08634, 2605.24774, 2605.00580, 2605.00105, 2605.29990 (KNOWN_PERL_ERRORS #387), 2605.15405
%          `h\,\partial f_A/\partial p_j`, 2605.00581 `\partial_{11}l(F(x),Y)f(x)`; the 57bw A/B XML (abg13, 3,003 papers)
%          applied a `\partial` to a greedy slash quotient in 370 formulas / 123 papers, a product of derivatives in 807 / 90,
%          another product in 916 / 94, a MulOp chain in 42 / 24
% oracle:  Perl 0 errors, its greedy `bigop` (MathGrammar:717, :605-618): `partial-differential@(V / partial-differential@(theta))`,
%          `T * partial-differential@(F / partial-differential@(T))`, `partial-differential@(Omega * open-closed-interval@(0, T))`;
%          user ruling 2026-09-29: ∂ takes one factor (`\partial_x u\cdot v` (∂_x u)·v) while ∇ keeps its greedy bare
%          argument (`\nabla u\cdot v` ∇@(u·v)); ∂F/∂T is a derivative; ∂Ω×(0,T] is ∂Ω times an interval
% cause:   `bigop_application = any_bigop term` took a DIFFOP too; divergence #374 (`diffop_application`,
%          `differential_operator_apply`, `leibniz_quotient`, `regroup_leibniz_numerator`)
\section{A differential operator takes one factor}
$\partial_x u\cdot v$ $\partial\Omega\times(0,T]$ $g^{\mu\nu}\partial_\mu\varphi\partial_\nu\varphi$ $\partial_x u\,\partial_y v$
$\nu\partial_x u\cdot\partial_x v$ $\partial_t u+\partial_x f(u)=0$ $\partial_\mu A^\mu$ $\partial_\mu\partial^\mu\phi$
$\partial_t\int_\Omega u\,dx$ $\int\partial_x u\,dx$ $\partial\sin x\cdot y$ $\partial\log x\cdot y$ $\nabla\partial_x u$
$\sin\partial_x u$ $\partial^\rho G(x-y)c(y)$ $\partial_{11}l(F(x),Y)f(x)$ $\partial f(x)$ $\nabla u\cdot v$
$\nabla\cdot u$ $\partial\Omega$ $\int_{\partial\Omega}u\,dS$ $\partial/\partial t$ $\partial_x+\partial_y$

% A Leibniz quotient is one derivative, whatever stands around it: the numerator is the last differential operator
% with the factors after it, the denominator the differentials its order asks for, and the rest multiplies the
% quotient (`leibniz_quotient`); a `\frac` numerator over a differential operator reads as the slash one does
% (`regroup_leibniz_numerator`, parser.rs `is_leibniz_numerator_arg`)
\section{A Leibniz quotient is one derivative}
$\partial F/\partial T$ $T\,\partial F/\partial T$ $h\,\partial f_{A}/\partial p_{j}$ $\partial^2 f/\partial x\partial y$
$\partial u/\partial x\,v$ $\partial\rho u/\partial t$ $\partial u/\partial x+\partial v/\partial y$ $\frac12\partial u/\partial x$
$\partial^n f/\partial x_1\cdots\partial x_n$ $\frac{\partial^2 u}{\partial x^2}$ $\frac{\partial\rho u}{\partial t}$
$\frac{\partial F}{\partial T}$

\end{document}
